% !TeX root = ../tfg.tex
% !TeX encoding = utf8
%
%*******************************************************
% Summary (800--1500 palabras, en inglés)
%*******************************************************

\selectlanguage{english}
\chapter{Summary}

\paragraph{Context and motivation.} Histopathological analysis of tissue samples is the clinical reference for diagnosing and grading prostate cancer. Stained with haematoxylin and eosin (H\&E) and digitised as whole-slide images, the samples are usually split into small patches that a deep neural network classifies one by one. This work uses a public dataset of \num{126109} prostate patches of $224 \times 224$ pixels annotated by pathologists into 14 histological classes, finer than the Gleason grades. A classifier trained on these classes always answers with one of them, even when the image comes from another organ, another hospital or a degraded acquisition, and it gives no reason for its decision. European regulation treats medical AI systems as high-risk and demands transparency, explainability, human oversight and robustness. A trustworthy system should therefore abstain when it faces inputs unlike its training data, and explain its decisions faithfully. Both requirements are the subject of this work: \emph{out-of-distribution} (OOD) detection and explainability, applied to the same prostate classifier.

\paragraph{Objectives and research questions.} The general objective is to study OOD detection and explainability methods that let a professional decide when to trust a histopathology classifier and understand what its decisions rely on, with a mathematical foundation for every method and a careful experimental evaluation. It is broken down into six specific objectives: to formalise the mathematical foundations; to train and compare five pretrained architectures; to design an OOD benchmark that separates semantic shift, change of centre, resolution and synthetic corruptions; to evaluate post hoc detectors on it and check whether the conclusions depend on the classifier; to reimplement and analyse STOOD-X, an explainable OOD detector; and to apply Layer-wise Relevance Propagation (LRP) and Concept Relevance Propagation (CRP) to the classifier. Seven research questions guide the work, from which architecture classifies best to whether relevance-based explanations are faithful enough to reveal biases.

\paragraph{Mathematical part.} The theoretical chapters formalise the tools used later. Supervised learning is presented through feed-forward and convolutional networks, residual connections, self-attention and the architectures compared, together with the cross-entropy loss, label smoothing and the AdamW optimiser. OOD detection is defined as a decision rule that thresholds a score, and its evaluation metrics are derived from the ROC curve; in particular, the area under the curve (AUROC) is shown to equal the Mann-Whitney statistic normalised by the product of the sample sizes, $U/(n_1 n_2)$, that is, the probability that a random OOD image receives a higher OOD score than a random in-distribution one. The post hoc detectors are described by family: those based on logits, those that modify activations, and those based on features and distances, including the Mahalanobis distance and STOOD-X, which runs a Mann-Whitney test between the distances of an image to its nearest training neighbours and those of the neighbours to their own. The inference chapter covers hypothesis testing, the Mann-Whitney test and its normal approximation with ties, McNemar's test for paired classifiers, the Friedman and Nemenyi tests for several methods over several datasets, quantile-quantile plots, bootstrap confidence intervals and conformal prediction, whose split version guarantees a bounded false-alarm rate under exchangeability. The last theoretical chapter presents LRP, its propagation rules and their analysis as a Markov chain, and CRP, which conditions relevance on internal concepts.

\paragraph{Computer science and experimental part.} Five ImageNet-pretrained architectures (ViT-B/16, ResNet-50, ConvNeXt-S, EfficientNetV2-S and Swin-S) were fine-tuned with fastai and timm on a GPU cluster managed with SLURM. Swin-S was then integrated into the OpenOOD benchmark to evaluate 13 configurations of 11 post hoc detectors and four configurations of STOOD-X. The OOD benchmark combines near-OOD colon tissue from NCT-CRC (tumour and normal mucosa); an internal control with colon tissues that resemble prostate classes (stroma, lymphocytes and smooth muscle); prostate tiles from other centres and at a much lower magnification (PANDA); a resolution control obtained by downscaling the in-distribution test images; nine synthetic corruptions common in digital pathology at three severity levels; and a far-OOD texture dataset (DTD). Every metric is reported with bootstrap confidence intervals, and the whole OOD evaluation was repeated with ConvNeXt-S. STOOD-X was reimplemented following its paper, its number of neighbours $K$ was selected with a leave-one-class-out validation that uses in-distribution data only, and its decision threshold was calibrated with conformal prediction. Finally, LRP and CRP were applied to ConvNeXt-S with zennit and zennit-crp, after choosing how to treat its LayerNorm layers, and were compared with gradient-based methods (Expected Gradients and Gradient $\times$ Input) in terms of faithfulness, concept ablation and RelMax concept atlases.

\paragraph{Main results.} Swin-S reaches an accuracy of 0.9688 [0.9658; 0.9717] on the 14 classes, but its advantage over ConvNeXt-S (+0.0020) is not significant (McNemar, $p = 0.19$); both models must be considered tied. When the classes are grouped into benign and malignant tissue, the accuracy of Swin-S rises to 99.33\,\%. Feature- and distance-based detectors consistently outperform logit-based ones and those that modify activations, which were designed for networks with ReLU and BatchNorm and transfer poorly to a Swin Transformer. The internal control shows that the detectors respond mainly to the acquisition conditions rather than to the tissue type: lymphocytes, a tissue also present in the prostate, are rejected almost as well as colon tissue. With $K = 144$, selected without any OOD data, STOOD-X is the best detector, with an AUROC of 0.991 and a false positive rate at 95\,\% true positive rate (FPR@95) of 0.026 on colon tissue, and a mean AUROC of 0.972 over eight OOD datasets. It is also the most sensitive detector to synthetic corruptions, although the saturation of its statistic limits its ability to grade their severity. Its $p$-values, however, are not calibrated: with the paper's value $K = 500$, 82\,\% of the in-distribution validation images obtain $p < 0.05$. Conformal calibration of the threshold keeps the false-alarm rate on the in-distribution test set compatible with the chosen level in all fifteen cases evaluated. Repeating the experiment with ConvNeXt-S changes several conclusions, such as the ranking of the detectors on near-OOD data. In explainability, the LRP variant used preserves relevance but is barely faithful: its maps hardly change when the model weights are randomised, and the channels it selects weigh little in the prediction. Gradient-based methods are more faithful, but their concept-conditioned maps are dominated by sampling noise. The analysis nevertheless revealed that selecting features by the sum of their absolute values favours a single high-magnitude channel that carries no class information, and that the border between tissue and white background recurs in the explanations. \revisar{depende de la conclusión del cap. 12/13}

\paragraph{Conclusions, limitations and future work.} Post hoc OOD detection is viable in this problem, and STOOD-X, with a number of neighbours suited to the class sizes, is the most effective of the detectors evaluated. However, changes in acquisition conditions are still hard to separate from the in-distribution data: detecting 95\,\% of those images requires rejecting too many prostate images for clinical use, and the conclusions depend on the classifier. The evaluation has several limitations. The dataset partition was made by patch, not by patient, so patches from the same slide appear in training and test, and the metrics measure generalisation within the same slides. The validation set was used for early stopping, so its exchangeability with the test set, which conformal guarantees require, is not ensured. Most of the evaluation uses a single model. Future work includes retraining with a patient-level split or validating on an external prostate dataset, studying the calibration of STOOD-X with $K = 144$, adapting LRP rules to attention blocks so that the Swin Transformer can be explained, and intervening on the suspected biased concepts to measure their effect on accuracy and OOD detection.

\medskip
\noindent\textbf{Keywords:} out-of-distribution detection, explainable artificial intelligence, digital histopathology, prostate cancer, STOOD-X, conformal prediction, Layer-wise Relevance Propagation, Concept Relevance Propagation.

% Al finalizar el resumen en inglés, volvemos a seleccionar el idioma español para el documento
\selectlanguage{spanish}
\endinput
